\documentclass[UTF8,11pt,a4paper]{ctexart}

\usepackage[a4paper,margin=2.4cm]{geometry}
\usepackage{amsmath,amssymb,amsthm,mathtools,bm}
\usepackage{booktabs}
\usepackage{enumitem}
\usepackage{graphicx}
\usepackage{hyperref}
\usepackage{xcolor}

\hypersetup{
  hidelinks
}

\setlist[itemize]{topsep=3pt,itemsep=2pt,parsep=0pt,partopsep=0pt}
\setlist[enumerate]{topsep=3pt,itemsep=2pt,parsep=0pt,partopsep=0pt}
\setlength{\parindent}{2em}
\setlength{\parskip}{0.35em}
\linespread{1.12}

\newtheorem{definition}{Definition}
\newtheorem{proposition}{Proposition}
\newtheorem{lemma}{Lemma}
\newtheorem{example}{Example}
\newtheorem{remark}{Remark}

\DeclareMathOperator*{\argmax}{arg\,max}
\DeclareMathOperator*{\argmin}{arg\,min}
\DeclareMathOperator*{\plim}{plim}

\newcommand{\E}{\mathbb{E}}
\newcommand{\Var}{\mathrm{Var}}
\newcommand{\Cov}{\mathrm{Cov}}
\newcommand{\Prob}{\mathbb{P}}
\newcommand{\R}{\mathbb{R}}
\newcommand{\dd}{\,\mathrm{d}}
\newcommand{\ind}{\mathbf{1}}
\newcommand{\pd}[2]{\frac{\partial #1}{\partial #2}}
\newcommand{\pdd}[2]{\frac{\partial^2 #1}{\partial #2^2}}
\newcommand{\given}{\,\middle|\,}

\title{ {{COURSE_TITLE_ZH}} \\ \large {{LECTURE_TITLE}} }
\author{}
\date{{{LECTURE_DATE}}}

\begin{document}

\maketitle

\section{本讲概览}

\begin{itemize}
{{SUMMARY_ITEMS}}
\end{itemize}

\section{核心内容}

% Use this section for the main lecture flow.
% Split into subsections when the logic clearly changes.

\subsection{问题背景}

% Summarize the motivating question in concise Chinese prose.

\subsection{模型设定或概念框架}

% State assumptions, objects, and notation here.

\section{推导与证明}

% Put core derivations here. Use align for multi-line math.

\section{经济学直觉与结论}

% Explain what the formulas mean and why the result matters.

\section{遗留问题}

\begin{itemize}
  \item 若有符号、板书或图形仍需核对，在这里列出。
\end{itemize}

\end{document}
